\documentclass[11pt]{article}
\usepackage[margin=25mm]{geometry}
\usepackage{fontspec}
\setmainfont{Latin Modern Roman}
\usepackage{amsmath,amssymb,mathtools}
\usepackage{unicode-math}
\setmathfont{Latin Modern Math}
\usepackage{xurl,hyperref}
\hypersetup{colorlinks=true,urlcolor=blue}
\setcounter{secnumdepth}{0}
\setlength{\parindent}{0pt}
\setlength{\parskip}{5pt}
\setlength{\emergencystretch}{3em}
\newcommand{\sref}[2]{\hyperlink{source:#1:#2}{[S#2]}}
\newcommand{\cataloguescope}[1]{\paragraph{Recorded target.}#1}
\newcommand{\R}{\mathbb R}
\newcommand{\C}{\mathbb C}
\newcommand{\N}{\mathbb N}
\newcommand{\Z}{\mathbb Z}
\newcommand{\Q}{\mathbb Q}
\newcommand{\F}{\mathbb F}
\newcommand{\eps}{\varepsilon}
\begin{document}
% UnsolvedMath ID 8000007; source catalogue code AMR-079-0007
\setcounter{section}{8000006}
\section{Non-self adjoint spectral problems}\label{8000007}
{\small UnsolvedMath ID 8000007 · Mathematical Physics}
\subsection{Definitions and mathematical statement}

\paragraph{Shared notation.}$\mathbb N=\{1,2,\ldots\}$; the standard number fields have their usual meanings. All additional conventions are stated locally.


\paragraph{Definitions and notation.}For a closed operator $A$, the spectrum $\sigma(A)$ consists of parameters where $A-\lambda$ has no bounded inverse. Its $\varepsilon$-pseudospectrum also includes parameters where $\|(A-\lambda)^{-1}\|>\varepsilon^{-1}$. Numerical range consists of $\langle Au,u\rangle$ for unit vectors in the domain. A semiclassical pseudomode is a unit vector with residual tending rapidly to zero as $h\to0$; it need not correspond to an actual eigenvalue.

\paragraph{Catalogue wording retained for traceability.}
Much is now known, theoretically and also numerically, about the spectrum of self-adjoint operators. This is in great contrast to the situation regarding non-self-adjoint operators, where the spectral theory is far more subtle and numerical schemes must overcome significant, inherent instabilities. The heart of the difficulties and instabilities lies in the following fact: in the self-adjoint case $A=A^*$, the resolvent satisfies $\|(A-\lambda)^{-1}\|=1/\operatorname{dist}(\lambda,\sigma(A))$ for $\lambda\notin\sigma(A)$, but in the non-self-adjoint case, this is no longer true, as we see already in the $2\times2$ case, $A=\begin{bmatrix}0&n\\0&0\end{bmatrix}$, $n\to\infty$. In recent years, a number of authors (e.g. Trefethan, Davies,...) have initiated a systematic approach to non-self-adjoint spectral problems, notions such as pseudospectrum have come into prominence, and other authors have conducted in depth studies of particular non-self-adjoint spectrum problems which arise in practice. For example, in analyzing the semi-classical limit of the focusing NLS equation, one must analyze the spectrum of the associated AKNS operator $T(h)$, as Planck's constant $h$ goes to zero. The operator is non-self-adjoint, and as $h$ goes to zero more and more eigenvalues, corresponding to solitons, emerge in the complex plane. It is of critical importance to the analysis of the semi-classical limits for NLS to determine where in the plane, and at what rate, the eigenvalues accumulate. The difficulty in doing this, theoretically and numerically, is illustrated by the following fact: the spectrum of $T(h)$ off the real axis is a discrete set, whereas the numerical range of $T(h)$ is an open subset of the plane. Nevertheless, every point $\lambda$ lying in the numerical range of $T(h)$ is an eigenvalue of $T(h)$ to all orders in $h$, $\|(T(h)-\lambda)u\|=O(|h|^k)$ for any $k\geq 1$, for some $u=u(h)$, $\|u\|=1$. In the language of Kruskal, the computation of the spectrum of $T(h)$ is a problem ``beyond all orders.'' Much has been done (Kamvissis-Miller-McLaughlin, Tovbis-Venakides-Zhou,...) in analyzing the semi-classical limit for NLS in special situations where the spectral problem can be solved explicitly. The spectral problem with general data, both for NLS and also other related non-self-adjoint problems, however, is far from understood, and poses a great challenge whose resolution is still only in the initial stages.
\paragraph{Scope and source qualification.}The source’s programme includes eigenvalue accumulation location and rate for general-data AKNS operators. Its informal assertion about pseudomodes throughout the numerical range is source context, not an added universal theorem for arbitrary operators. The source’s distance wording is corrected to reciprocal distance: for self-adjoint A the resolvent norm is the reciprocal of the distance to its spectrum, at resolvent parameters.
\subsection{Short English statement}
Understand and compute non-self-adjoint spectra, especially the semiclassical scattering spectrum of focusing nonlinear Schrödinger equations.
\subsection{Research attempt}
\paragraph{Model, date and investigated scope.}
OpenAI GPT-6 Astra (configured model \texttt{gpt-6-astra}), ultra reasoning; generated 11 October 2026. Verified exact pseudospectral and perturbation formulas and an all-orders pseudomode with no nearby eigenvalue. General AKNS eigenvalue accumulation and certified computation remain unresolved.

\paragraph{Sources and literature check.}
Read the supplied complete definition, catalogue wording, scope qualifications and source list for this canonical ID. The original-source citations are inherited as supplied references; no fresh primary-source or novelty search was performed for this attempt. No current open-status claim is inferred from this calculation.

\paragraph{An exact instability calculation.}For \(n>0\), let \(A_n=\begin{pmatrix}0&n\\0&0\end{pmatrix}\). At \(z\ne0\),
\[(A_n-zI)^{-1}=\begin{pmatrix}-z^{-1}&-nz^{-2}\\0&-z^{-1}\end{pmatrix},\qquad
\|(A_n-zI)^{-1}\|\ge n/|z|^2.
\]
Thus the disk \(0<|z|<\sqrt{n\varepsilon}\) lies in its \(\varepsilon\)-pseudospectrum, although the spectrum consists only of zero. Perturbing the lower-left entry by \(\delta\) gives characteristic equation \(z^2=n\delta\). A perturbation of norm \(|\delta|\) can therefore create eigenvalues at distance \(\sqrt{n|\delta|}\).
\paragraph{Residuals beyond all algebraic orders.}Choose \(n(h)=e^{1/h}\) for \(h>0\) and fix \(z\ne0\). The normalized vector proportional to \((1,z/n(h))\) has residual under \(A_{n(h)}-zI\) of norm at most \(|z|^2e^{-1/h}\). This is smaller than every positive power of \(h\), while \(z\) remains a fixed positive distance from the exact spectrum. To see the last asymptotic directly, \(e^{-1/h}/h^m\to0\) after setting \(s=1/h\), since \(e^s\) dominates its polynomial Taylor term of degree \(m+1\).
\paragraph{What the example rules out.}An all-orders residual estimate alone cannot certify an AKNS eigenvalue or its limiting distribution. A numerical certificate must control a resolvent or characteristic function on a contour in a norm adapted to the operator, with errors below the sensitivity scale. The two-by-two example is not an AKNS operator, and its exponentially growing coefficient is not claimed to model a fixed smooth potential. It only disproves the logical inference from tiny residual to nearby spectrum for arbitrary nonnormal families. The source's general-data semiclassical accumulation location and rate remain unproved, as does a stable numerical method for them.

\paragraph{Verification and adversarial audit.}
Checked inverse signs, norm lower bound, perturbation norm, normalized residual and explicit distinction between the matrix example and an AKNS family. This was an LLM adversarial audit, not independent human verification; no theorem prover was used. Prompts A then B applied to this ID: Checked inverse signs, norm lower bound, perturbation norm, normalized residual and explicit distinction between the matrix example and an AKNS family. Scope remains limited: Obtain operator-specific resolvent or characteristic-function control that separates actual AKNS eigenvalues from pseudomodes.

Coordinator checked the Jordan-block resolvent, perturbation eigenvalues and exponentially small pseudomode residual. The example is not presented as an AKNS operator. This additional check was model-assisted, not independent human verification.

\paragraph{Tools, process and contributions.}
Python 3 and local shell used to read supplied files and serialize this manuscript; mathematics derived symbolically in the text, with no numerical experiment or formal proof assistant. Contribution prompt A was followed by prompt B. GPT-6 Astra with ultra reasoning drafted the research and performed its stated self-audit; any additional coordinator contribution is identified in the verification and audit. A GPT-6-based Codex coordinator assembled the manuscript and checked repository integration; its exact variant and reasoning configuration were not exposed. No human mathematical contribution, independent human proof check or formal proof certification is supplied. Complete derivations and the consulted source references are included here; private reasoning transcripts are not published.

\paragraph{Final status and first remaining gap.}
\textbf{UNRESOLVED}. Obtain operator-specific resolvent or characteristic-function control that separates actual AKNS eigenvalues from pseudomodes.

\section{Completion Estimate}
\noindent\textbf{COMPLETION: 0\%}
The calculation is an elementary benchmark or conditional reduction; it establishes no claimed new portion of the complete source target.

\subsection{Sources}
\begin{itemize}
\item[\hypertarget{source:8000007:1}{S1}] ULAM.AI and the UnsolvedMath Contributors, UnsolvedMath, record 8000007 (AMR-079-0007). Catalogue curation under CC BY 4.0; definitions and English exposition added locally.
\url{https://huggingface.co/datasets/ulamai/UnsolvedMath}.
\item[\hypertarget{source:8000007:2}{S2}] Percy Deift, Some open problems in random matrix theory and the theory of integrable systems (2007), arXiv:0712.0849, Problem 7.
\url{https://arxiv.org/abs/0712.0849}.
\end{itemize}
\label{end:8000007}
\end{document}
